\chapter{Diachrony of directional extensions} \label{ch:diachrony} 

This chapter compiles the available evidence for how directional extensions have developed historically in Chadic languages. Historical reconstruction has been a focus for many early scholars of Chadic languages, but beyond the general internal structure of the family proposed by \citet{newmanChadicClassificationReconstructions1977}, there has only been limited progress in developing an overall picture of the history of these languages. In the assessment of \citet[343--344]{guldemannHistoricalLinguisticsGenealogical2018}, ``the reconstructions [of lexical forms]... are not established transparently in a bottom-up procedure within a clear phylogenetic structure, and only a limited number of them involve the level of the proto-language'' and in regard to grammaticalized morphemes ``hardly any domain has received such a depth of research as to produce a concrete set of morphological proto-forms that are based either on subgroup reconstructions or, given the size of the family, at least on a representative language sample.''

\section{Previous reconstructions of directional extensions}

In regard to directional extensions, there have been a few proposals. \citet[275]{newmanChadicExtensionsPredative1977} claims that ventive extensions ``can be reconstructed for Proto-Chadic with a high degree of certainty.'' Two forms are reconstructed, each based on data from three West Chadic languages and three Biu-Mandara languages. Newman proposes that neither of these forms had a ventive function in Proto-Chadic. The proto-form $^\ast$\textit{(a)wa} is proposed to have been a distal demonstrative and the proto-form  $^\ast$\textit{in} is proposed to have been a benefactive marker.\footnote{Both distal location and benefactive meanings are co-expressed with ventive directionals in Chadic languages (Chapter \ref{ch:functions}, Sections \ref{sec:ventloc} and \ref{sec:beneficiary}), but it is unclear why the path of grammaticalization should be from location and beneficiary to ventive direction rather than the other way around. Since reflexes of Newman's ``distal'' $^\ast$\textit{(a)wa} can also express benefactive meaning, Newman suggests that the path of grammaticalization actually goes both ways. Not only can a benefactive marker become a ventive marker, but (a distal marker that becomes) a ventive marker can also become a benefactive marker.} It is left unexplained why many West Chadic languages have reflexes of both of these forms as aspect-related allomorphs. 

%\footnote{``The original meaning of the Distant extension would have been to indicate that the action of the verb was done ``there'' at some distance from the speaker... was probably entirely locative in nature'' \citep[276]{newmanChadicExtensionsPredative1977}.} 
%\footnote{``The primary meaning of this extension was to indicate that the action of a verb was destined for, done for the benefit of, or otherwise affected or pertained to someone'' \citep[281]{newmanChadicExtensionsPredative1977}.} 

\citet{frajzyngierVentiveCentrifugalChadic1987} takes a step forward by increasing the sample to 24 languages with ventive (and itive) directional extensions. Frajzyngier takes a more conservative approach by not proposing Proto-Chadic forms. For West Chadic ventive forms, the aspect-related allomorphy is explained by positing that these verbal extensions are derived from verbs in a directional multiverb construction. Frajzyngier notes that the forms of ventive and itive extensions in Biu-Mandara do not match the West Chadic forms and suggests that there may have been a variety of sources for directional extensions in Biu-Mandara in addition to the possible grammaticalization of verbs. The particular reconstructions proposed do not hold up to further scrutiny,\footnote{For example, \citet[35]{frajzyngierVentiveCentrifugalChadic1987} gives possible verbal sources of three itive directional extensions, but \hyperref[marg]{Margi} \textit{-ba} is an outward extension, \hyperref[musg]{Musgu} (Munjuk) \textit{-fu} was not found in other sources, and \hyperref[lagw]{Lagwan} \textit{-li} is an inward extension. \citet{heineWorldLexiconGrammaticalization2002} unfortunately refer to Frajzyngier's Lagwan data as an example of grammaticalization from a verb `go' to an itive directional marker. They also cite an example Frajzyngier gives from Hwana (hwan1240) but this data is from unpublished fieldnotes and cannot be verified.} but the key contribution of this paper is to observe that widespread use of ventive directional extensions does not necessarily entail a common etymological source, especially given the absence of cognate forms across branches of Chadic languages. 

\citet[306--318]{schuhChadicCornucopia2017} also recognizes the differences between directional extensions in each branch of Chadic languages. Three consonants are suggested as proto-forms of ventive directional extensions in West Chadic languages:  $^\ast$\textit{-w}, $^\ast$\textit{-n}, and $^\ast$\textit{-t}. Schuh tentatively suggests that these proto-forms may have shifted to other functions in Biu-Mandara languages, but largely concludes that Biu-Mandara languages ``share an areal feature of having developed complex systems of directional verbal extensions through reinterpretation and grammaticalization of verbs, nouns, adverbs, and prepositions'' \citep[316]{schuhChadicCornucopia2017}.\footnote{Schuh excludes the Kotoko-Buduma and Musguic languages from this description as part of his proposed ``Central Chadic B'' group.} 

\begin{figure}[t]
	\begin{forest}
		forked edges, for tree={grow=0,edge=thick, s sep=-4pt, anchor=base west,  font=\strut\footnotesize\sffamily, grow'=east} % s sep=0pt,
		[West Chadic  
		[A 
		[A1 [Hausa \textit{ō},fill={white!50!yellow}, tier=language] ]
		[A2 
		[Boleic 
		[Karekare \textit{(n)ee, tu} ,fill={white!50!yellow}, tier=language]
		[Beele \textit{u},fill={white!50!yellow}, tier=language]
		[Bole \textit{n, tV, aakoo},fill={white!50!yellow}, tier=language]
		[Geruma \textit{ìŋ},fill={white!50!yellow}, tier=language]
		[Giiwo \textit{(í)n, o},fill={white!50!yellow}, tier=language]
		[Ngamo \textit{no, n, t},fill={white!50!yellow}, tier=language]
		[Bure \textit{ín},fill={white!50!yellow}, tier=language]
		[Gera /L tone/, tier=language]
		[\textit{None} {(Nyam, Galambu)}, tier=language]
		]
		[Tangalic 
		[Kanakuru \textit{tə(ru), bo},fill={white!50!yellow}, tier=language]
		[Kushi \textit{na, ru},fill={white!50!yellow}, tier=language]
		[Pero \textit{na, tu},fill={white!50!yellow}, tier=language]
		[Tangale \textit{tu, na},fill={white!50!yellow}, tier=language]
		]
		]
		[A3 
		[\textit{None} {(Goemai, Mupun, Ngas)}]
		]
		[A4 
		[\textit{None} {(Fyer, Kulere, Ron-Daffo)}]
		[Ron-Scha \textit{ó} ,fill={white!50!yellow}]
		]
		] 
		[B 
		[B1 
		[Bade \textit{àawo, ìina} ,fill={white!50!yellow}]
		[Ngizim \textit{ay, iina} ,fill={white!50!yellow}]
		]
		[B2 
		[\textit{None} ({Pa'a, Miya})]
		]
		[B3 
		[\textit{None} ({Guruntum, Dass, Zaar})]
		]
		] 
		]
	\end{forest}
	\caption{Ventive forms in West Chadic} \label{fig:westventive}
\end{figure}

\section{Ventive}

Following previous work which has not established links between ventive directional extensions across the branches of Chadic languages, the history of directional extensions is discussed separately here for each branch. 

\subsection{West Chadic}

Nearly all ventive forms in West Chadic languages resemble at least one of the three proposed reconstructions, $^\ast$\textit{-w}, $^\ast$\textit{-n}, and $^\ast$\textit{-t}, as marked in yellow in Figure \ref{fig:westventive}. One outlier is \hyperref[gera]{Gera} where ventive direction is expressed by a low tone. \citet[309]{schuhChadicCornucopia2017} does not identify any source of ventive morphemes in West Chadic languages and assumes that languages without a ventive extension have lost this marker, or its function and form have changed beyond recognition. To support this hypothesis, more research is needed into whether West Chadic languages without ventive extensions show evidence of erstwhile ventive extensions that have shifted into another function. The alternative proposal, suggested by \citet{frajzyngierVentiveCentrifugalChadic1987}, is that ventive directional extensions developed from directional multiverb constructions. This leaves open the possibility that this innovation occurred after West Chadic languages split off and did not necessarily occur in all West Chadic languages (e.g., A3, B2, B3). The path of grammaticalization from ventive verb to ventive directional marker is attested in various languages around the world \citep[85--86]{heineWorldLexiconGrammaticalization2002}. However, the reconstruction of this ventive verb root is much less certain. Frajzyngier's proposal assumes a Proto-Chadic verb $^\ast$\textit{wat} as the source, but this is a very speculative reconstruction. Many West Chadic languages (and Biu-Mandara) languages have a motion verb beginning with \textit{w} or \textit{nd}/\textit{d}/\textit{dz}; however, these are most often itive or general motion verb roots (not ventive).  

\subsection{Biu-Mandara and Masa}

\begin{figure}[t]
	\begin{forest}
		forked edges, for tree={grow=0,edge=thick, s sep=-2.5pt, anchor=base west,  font=\strut\footnotesize\sffamily, grow'=east} % s sep=0pt,
		[Biu-Mandara
		[North 
		[Gidar [\textit{None} (Gidar)] ]
		[Higic 
		[Bana \textit{ke} ,fill={white!50!lightgray}]
		[Psikye \textit{ke}? ,fill={white!50!lightgray}]
		]
		[Kotoko-Buduma 
		[\textit{None} ({Buduma, Lagwan, Makary Kotoko})]
		]
		[Lamang-Hdi [\textit{None} (Lamang) ] ]
		[Ma.-Ma.-Mo.	 
		[Bura-Marghi
		[\textit{None} ({Bura, Cibak, Margi, Huba}), tier=language]
		]
		[Mandaraic
		[Dghwede	\textit{àwá} ,fill={white!50!cyan}, tier=language]
		[Matal	\textit{awâŋ} ,fill={white!50!cyan}, tier=language]
		[Parkwa	\textit{ətsa, atsa}, tier=language]
		[\textit{None} ({Glavda, Malgwa, Wandala}), tier=language]
		]
		[Mofuic
		[Merey	\textit{aw} ,fill={white!50!cyan}, tier=language]
		[Zulgo-Gemzek	\textit{ara, aw, ya} ,fill={white!50!cyan}, tier=language]
		[Mofu-Gudur	\textit{wa} ,fill={white!50!cyan}, tier=language]
		[Mada	\textit{re}, tier=language]
		[Moloko	\textit{alaj}, tier=language]
		[Muyang	\textit{biyu}, tier=language]
		[Wuzlam	\textit{ārá}, tier=language]
		]
		]
		[Maroua [North Giziga	\textit{áwà, óò, o} ,fill={white!50!cyan}] ]
		[Musguic 
		[Musgu	\textit{sí} ,fill={white!20!pink}]
		[Mbara	\textit{sì} ,fill={white!20!pink}]
		]
		]
		]
	\end{forest}
	\caption{Ventive forms in Biu-Mandara North} \label{fig:ventivenorthbm}
\end{figure}

It is not yet possible to reconstruct a ventive form for Biu-Mandara languages. The forms of ventive directional extensions in North Biu-Mandara languages are shown in Figure \ref{fig:ventivenorthbm}. There are several forms similar to West Chadic $^\ast$\textit{awa}, primarily found in Margi-Mandara-Mofu languages. These are marked in blue. The Higic languages have a different form, \textit{ke} (marked in grey), and the Musguic languages have a distinct form \textit{si} (marked in pink). This last form is notably similar to ventive directional verbs in Biu-Mandara languages which often have a single consonant \textit{s} as the verb root. There are also a number of forms that are not clearly cognate with any others. While the widespread use of ventive forms suggest that a ventive of some form was likely present in Proto-North-Biu-Mandara, there is not enough evidence available to construct its archaic form. 

The most common form of the ventive directional extension in Biu-Mandara South language is similar to \textit{aha}, as marked in blue in Figure \ref{fig:ventivesouthbm}. These forms are clearly cognate in Dabaic languages, and possibly in Bataic as well. It is tempting to posit a correspondence between \textit{w} and \textit{h} that would link the Dabiac ventive extensions to those in Biu-Mandara North, but \citet{wolffLexicalReconstructionCentral2023} does not find any evidence of such a reflex. A variant featuring \textit{k} appears in Matakam languages, marked in gray. Further information would be needed to decide whether any form can be reconstructed for all South Biu-Mandara languages. 

\begin{figure}
	\begin{forest}
		forked edges, for tree={grow=0,edge=thick, s sep=-3pt, anchor=base west,  font=\strut\footnotesize\sffamily, grow'=east} % s sep=0pt,
		[Biu-Mandara
		[Hurza 
		[Mbuko	\textit{ay} ,fill={white!50!yellow}]
		[Vamé	\textit{aya, ki} ,fill={white!50!yellow}]
		]
		[North ]
		[South 
		[Bataic [Bacama	\textit{á} ,fill={white!50!cyan}, tier=language]
		[Gude	\textit{ʸà} ,fill={white!50!cyan}, tier=language]
		]
		[Dabaic
		[Buwal\textit{	ā, xā} ,fill={white!50!cyan}, tier=language]
		[Gavar	\textit{(h)à} ,fill={white!50!cyan}, tier=language]
		[Daba	\textit{a, (á)ha} ,fill={white!50!cyan}, tier=language]
		[Mazagway	\textit{aha} ,fill={white!50!cyan}, tier=language]
		[Mbudum	\textit{áhà} ,fill={white!50!cyan}, tier=language]
		[Mina	\textit{á(hà)} ,fill={white!50!cyan}, tier=language]
		]
		[Matakam
		[Cuvok	\textit{ék} ,fill={white!50!lightgray}, tier=language]
		[Mafa	\textit{ká(ɗá)} ,fill={white!50!lightgray}, tier=language]
		]
		[Sukur [\textit{None} (Sakun)] , tier=language]
		[Teraic	[\textit{None} {(Ga'anda, Tera)}] , tier=language]
		]
		]
	\end{forest}
	\caption{Ventive forms in Biu-Mandara South and Hurza}  \label{fig:ventivesouthbm}
\end{figure}

The Hurza languages, marked in yellow in Figure \ref{fig:ventivesouthbm}, have their own form of the ventive directional extension. \citet[11]{gravinaVerbPhraseMbuko2001} points to an adverbial \textit{àhày} `to here -- toward deictic centre' as the source of the \hyperref[mbuk]{Mbuko} ventive suffix \textit{-ay}. The Hurza ventive extensions are similar in form to two of the three Masa languages with a ventive directional extension (\textit{áy} in \hyperref[muse]{Musey} and \textit{ey} in \hyperref[masa]{Masana}) while a different form \textit{fi} is found in the other Masa language, \hyperref[herd]{Herde}. The sporadic distribution of potential cognates within Biu-Mandara and between Biu-Mandara and West Chadic languages makes it unclear if these are distantly related forms or coincidental similarities.

\subsection{East}

The three East Chadic languages with a ventive directional extension all have different forms from each other. The East A language, \hyperref[lele]{Lele}, has the form \textit{yè} which could potentially be explained through contact with Masa languages. In East B, the \hyperref[kera]{Kera} form \textit{dà} and \hyperref[soko]{Sokoro} \textit{ti} both feature coronals stops, and so are not dissimilar to some of the West Chadic forms. However, this is not enough data to justify re-constructing a Proto-Chadic ventive extension.\footnote{Compare, for example, ventive directional markers in Berber languages which also tend to have a coronal stop. These are not considered cognates with Chadic ventive extensions, and likely have a more recent origin developed from demonstratives \citep{mettouchiGrammaticalizationDirectionalClitics2011}.} 

In at least three East Chadic languages, the verb `come' appears to be formed of a verb root `go' plus another morpheme: \textit{ɗe} `go' and \textit{ɗéːna} `come' in \hyperref[kwan]{Kwang}, \textit{há} `go' and \textit{háde} `come' in \hyperref[somr]{Somrai}, \textit{è} `go' and \textit{èjè} `come' \hyperref[lele]{Lele}. In Kwang and Somrai, there is no evidence that the syllable distinguishing `go' and `come' is a suffix, as it does not appear with other verbs. In Lele, the form is used as a ventive suffix elsewhere, while in the form \textit{èjè} `come' there is evidence that it is part of the verb root. Are the forms in Kwang and Somrai evidence of the development of ventive directional extensions in East Chadic, or of their loss which has left behind only fossilized traces of erstwhile verbal morphology?

\section{Itive}

\begin{figure}[b]
\fittable{
	\begin{forest}
		forked edges, for tree={grow=0,edge=thick, s sep=-3pt, anchor=base west,  font=\strut\footnotesize\sffamily, grow'=east} % s sep=0pt,
		[Biu-Mandara
		[North 
		[Gidar [\textit{None} (Gidar)] ]
		[Higic 
		[Bana \textit{ghe} ,fill={white!50!yellow}]
		[Psikye \textit{gu} ,fill={white!50!yellow}]
		]
		[Kotoko-Buduma 
		[\textit{None} ({Buduma, Lagwan, Makary Kotoko})]
		]
		[Lamang-Hdi [\textit{None} (Lamang) ] ]
		[Ma.-Ma.-Mo. 
		[Bura-Marghi
		[Bura \textit{mta} ,fill={white!50!cyan} ]
		[Cibak \textit{nta} ? ,fill={white!50!cyan}]
		[Margi \textit{nà} ,fill={white!50!cyan}]
		[Huba \textit{nà},fill={white!50!cyan} ]
		]
		[Mandaraic
		[Dghwede	\textit{dá} ,fill={white!50!lightgray}]
		[Glavda \textit{d(a)} ,fill={white!50!lightgray} ]
		[\textit{None} {(Parkwa, Malgwa, Wandala, Matal)}]
		]
		[Mofuic
		[Zulgo-Gemzek	\textcolor{red}{\textit{aha}} ]
		[Moloko	\textcolor{red}{\textit{ala}}]
		[Mada	\textit{ro},fill={white!50!lime}]
		[Wuzlam	\textit{ārāy, ə̄rə̄},fill={white!50!lime}]
		[\textit{None} ({Merey, Mofu-Gudur, Muyang})]
		]
		]
		[Maroua [\textit{None} (North Giziga)] ]
		[Musguic 
		[Musgu	\textcolor{red}{\textit{di}} ]
		[\textit{None} (Mbara) ]
		]
		]
%		[Hurza 
%		[\textit{None} {(Mbuko, Vamé)} ]
%		]
%		[South 
%		[Bataic [\textit{None} ({Bacama, Gude})	] ]
%		[Dabaic
%		[Buwal \textit{āzà} ,fill={white!20!pink}]
%		[Daba	\textit{ay} ,fill={white!20!pink}]
%		[\textit{None} {(Gavar, Mazagway, Mbudum, Mina)} ]
%		]
%		[Matakam
%		[Cuvok	\textcolor{red}{\textit{áɗ}}]
%		[\textit{None} (Mafa)	]
%		]
%		[Sukur [\textit{None} (Sakun)] ]
%		[Teraic	
%		[Ga'anda \textcolor{red}{\textit{kaɗe}} ] 
%		[\textit{None} (Tera)  ] ]
%		]
		]
	\end{forest}
	}
	\caption{Itive forms in Biu-Mandara North} \label{fig:itivebm}
\end{figure}

Itive directional extensions in Biu-Mandara languages have diverse forms, as shown in Figure \ref{fig:itivebm} and \ref{fig:itivesouthbm}. This presents a challenge for any effort to reconstruct such a marker for Proto-Biu-Mandara. There are plausible links to itive directional verb roots suggesting that at least some itive directional extensions may have formed in the context of multiverb constructions, as suggested by \citet{frajzyngierVentiveCentrifugalChadic1987}. The form \textit{-da} in \hyperref[dghw]{Dghwede} and \hyperref[glav]{Glavda} matches the itive verb root \textit{d-} found across all Mandaraic languages, but it is also similar to a preposition \textit{dà} `to(ward)' in Dghwede. Itive directionals are most commonly found within the Margi-Mandara-Mofuic subgroup, and within these languages there is a possibility that the coronal consonants seen in itive directional extensions are historically related.

There may be other historical links that have been obscured by semantic shift. The form of the itive directional extension \textit{-álâ} in \hyperref[molo]{Moloko} matches the form of a telicity-related aspectual suffix \textit{álâ} in Matal and a suffix \textit{-al} in \hyperref[glav]{Glavda} which often expresses the notion of separation. \citet[]{thomasGrammarSakunSukur2014} describes a ``centrifugal'' suffix \textit{-rá} in \hyperref[saku]{Sakun} which does not have an itive directional meaning in the examples given but could be historically related to itive directionals in other languages. In order to convincingly reconstruct a scenario in which directional extensions underwent a semantic shift, there would first need to be a comparative overview of verbal tense-aspect morphology in these languages.


\begin{figure}[t]
	\begin{forest}
		forked edges, for tree={grow=0,edge=thick, s sep=-3pt, anchor=base west,  font=\strut\footnotesize\sffamily, grow'=east} % s sep=0pt,
		[Biu-Mandara
		[Hurza 
		[\textit{None} {(Mbuko, Vamé)} ]
		]
		[North ]
		[South 
		[Bataic [\textit{None} ({Bacama, Gude})	] ]
		[Dabaic
		[Buwal \textit{āzà} ,fill={white!50!cyan}]
		[Daba	\textit{ay} ,fill={white!50!cyan}]
		[\textit{None} {(Gavar, Mazagway, Mbudum, Mina)} ]
		]
		[Matakam
		[Cuvok	\textit{áɗ} ,fill={white!50!lightgray}]
		[\textit{None} (Mafa)	]
		]
		[Sukur [\textit{None} (Sakun)] ]
		[Teraic	
		[Ga'anda \textit{kaɗe} ] 
		[Tera \textit{ɓara}] ]
		]
		]
	\end{forest}
	\caption{Itive forms in Biu-Mandara South and Hurza} \label{fig:itivesouthbm}
\end{figure}

The sole example of an itive directional extensions in East Chadic is found in \hyperref[kera]{Kera}. Notably, there is also a postposition of the same form. This suggests that the itive directional extension may be a recent innovation derived from the postposition. 

\section{Boundary-crossing}

There are some clusters of cognates for boundary-crossing directional extensions in Biu-Mandara languages, but it is not the case that one form can be identified as the etymological source of the majority of these directional extensions. Rather, it appears more likely that inward and outward directional extensions have developed independently several times, most often from adpositions, but perhaps also from verbs. 

\begin{table}[b]
	\caption{Forms of inward extensions (Biu-Mandara)}
	\label{tab:intolang}
	\begin{tabular}{lllll}
		\lsptoprule
		&     Forms    &  Subbranch & Group & Subgroup        \\
		\midrule
		Psikye	&	\textit{mbe}	&	North	&	Higic	&	\\
		Buduma	&	(\textit{a})\textit{l}	&	North	&	Kotoko-Buduma	&	\\
		Lagwan	&	\textit{li}	&	North	&	Kotoko-Buduma	& \\
		Lamang	&	\textit{ŋ}	&	North	&	Lamang-Hdi	& \\
		Bura	&	\textit{wa}, \textit{kwa}, \textit{u}	&	North	&	M.-M.-M.	& Bura-Marghi \\
		Margi	&	\textit{wá}	&	North	&	M.-M.-M.	& Bura-Marghi \\
		Dghwede	&	\textit{mè}	&	North	&	M.-M.-M.	& Mandaraic \\
		Parkwa	&	\textit{əkwa}, \textit{akwa}	&	North	&	M.-M.-M.	& Mandaraic \\
		Glavda	&	\textit{m}	&	North	&	M.-M.-M.	& Mandaraic \\
		Malgwa	&	\textit{ə́m}	&	North	&	M.-M.-M.	& Mandaraic \\
		Wandala	&	\textit{m}	&	North	&	M.-M.-M.	& Mandaraic \\
		Matal	&	\textit{ábà}, \textit{águ}	&	North	&	M.-M.-M.	& Mandaraic \\
		Mada	&	\textit{va}, \textit{a}	&	North	&	M.-M.-M.	& Mofuic \\
		Muyang 	&	\textit{iyu}	&	North	&	M.-M.-M.	& Mofuic \\
		Gude	&	\textit{gərə}	&	South	&	Bataic	& \\ 
		Sakun	&	\textit{ɣə}	&	South	&	Sukur	& \\
		\lspbottomrule
	\end{tabular}
\end{table}

\begin{table}[b]
	\caption{Forms of outward extensions (Biu-Mandara)}
	\label{tab:outlang}
	\begin{tabular}{lllll}
		\lsptoprule
		&     Forms    &  Subbranch & Group & Subgroup        \\
		\midrule
		Psikye 	&	(\textit{aka})\textit{ve}	&	North	&	Higic	&	\\	
		Lamang	&	\textit{b}(\textit{i})	&	North	&	Lamang-Hdi	&	\\	
		Bura	&	\textit{bila}, \textit{ɓəla}	&	North	&	M.-M.-M.	&	Bura-Marghi	\\
		Cibak	&	\textit{ba}	&	North	&	M.-M.-M.	&	Bura-Marghi	\\
		Margi	&	\textit{bá}	&	North	&	M.-M.-M.	&	Bura-Marghi	\\
		Huba	&	\textit{bìyà}	&	North	&	M.-M.-M.	&	Bura-Marghi	\\
		Dghwede	&	\textit{àwílè}	&	North	&	M.-M.-M.	&	Mandaraic	\\
		Parkwa	&	\textit{edá}, \textit{əlá}, \textit{adá}	&	North	&	M.-M.-M.	&	Mandaraic	\\
		Malgwa	&	\textit{sə́}	&	North	&	M.-M.-M.	&	Mandaraic	\\
		Wandala	&	\textit{s}	&	North	&	M.-M.-M.	&	Mandaraic	\\
		Mada	&	\textit{ra}	&	North	&	M.-M.-M.	&	Mofuic	\\
		Muyang	&	\textit{aya}	&	North	&	M.-M.-M.	&	Mofuic	\\
		Gude	&	\textit{gi}	&	South	&	Bataic	&		\\
		Sakun	&	\textit{va}	&	South	&	Sukur	&		\\
		Ga'anda	&	\textit{kaɗe}	&	South	&	Sukur	&		\\
		\lspbottomrule
	\end{tabular}
\end{table}

The forms of inward extension are shown in Table \ref{tab:intolang}. There are at least three clusters of likely cognate forms. The two Kotoko-Buduma languages have forms with a lateral: \textit{(a)l} (Buduma) and \textit{li} (Lagwan). These forms could be related to a \hyperref[marg]{Margi} verb \textit{li} which expresses inward direction. However, two other languages in the Bura-Margi cluster have the form \textit{li} as an itive directional verb root. This creates some uncertainty about the links between these forms, but also suggests possible diachronic links between itive and inward direction.

\hspace*{-1pt}The two Bura-Marghi languages have similar forms, \textit{wa} (Margi) and \textit{kwa} (Bura). This form also appears in the two other Bura-Marghi languages. In \hyperref[huba]{Huba}, the verb \textit{gwà(kwà)} `enter' is attested with and without the final syllable \textit{kwa} and in \hyperref[ciba]{Cibak} the verb \textit{lə̀kwá} `enter' is analyzed as a contraction of \textit{li} `go' and the preposition \textit{akwa} `into' \citep[136]{hoffmannZurSpracheCibak1955}. This syllable is not analyzed as a suffix in this language as it does not seem to appear with other verb roots. These forms appear to be cognates with \textit{akwə} in Parkwa. Other inward directional extensions with velar consonants could also be related, such as \textit{gərə} in Gude and \textit{ɣə} in Sakun.

The third cluster of cognates are the bilabial nasals founds in four of five Mandaraic languages. The more distantly related language Psikye has a form \textit{mbe} which also has a bilabial nasal. In Psikye, \citep[118]{smithKapsikiLanguage1969} points out that \textit{-mbe} is identical to the preposition \textit{mbe} `in' and is only attested with two verbs. The same form occurs in the closely related language \hyperref[bana]{Bana} but since it is only found with one verb, \textit{dzə̀mbə̀} `enter', it is not analyzed as a suffix. The Psikye and Bana forms may have more recently developed independently of the Mandaric inward directional extensions. Recent development from a preposition is also found in Sakun where the  inward suffix \textit{-ɣə} is similar to the preposition \textit{ɣi} ‘in’ \citep[273]{thomasGrammarSakunSukur2014}. The Mada inward directional extension \textit{-va} may have also recently developed from an adposition. The closely related language \hyperref[molo]{Moloko} has an ``adpositional enclitic'' \textit{ava} `in'.

Table \ref{tab:outlang} shows the forms of outward extensions. The Cibak and Margi forms are nearly identical, \textit{ba}.  \citet[125]{hoffmannGrammarMargiLanguage1963} suggests the Margi suffix (and presumably Cibak also) is derived from an outward directional verb \textit{bá}. Outward extensions in the other two languages in the Bura-Marghi subgroup (as well as Lamang) also begin with a bilabial voiced stop. Elsewhere, it is less clear what the sources of outward directional extensions may have been, with the possible exception of Psikye where the outward suffix is identical to a preposition \textit{ve} `to'. 

\section{Vertical}

The forms of upward extensions in Chadic languages are shown in Table \ref{tab:uplang} and downward extensions are shown in Table \ref{tab:downlang}. In both cases, the forms are diverse and the sporadic similarities cannot immediately be reconstructed into proto-forms. A cluster of possible cognate forms can be identified in the downward directional extensions of three languages of the Bura-Marghi subgroup: \textit{hi} (Bura), \textit{ía} (Margi) and \textit{yà} (Huba). For \hyperref[bura]{Bura}, \citet[282]{hoffmannUntersuchungenZurStruktur1955} compares the downward directional extension to a noun \textit{hi} `earth, soil' -- a path of grammaticalization attested in Oceanic languages and elsewhere \citep[121--122]{heineWorldLexiconGrammaticalization2002}. The two Biu-Mandara South languages, Sakun and Ga'anda, spoken in relatively close areas along the Nigeria-Cameroon border, both have a form \textit{xa} for the downward extension. In Sakun, the downward directional extension and the upwward directional extension are identical in form to prepositions with similar meanings \citep[273--274]{thomasGrammarSakunSukur2014}. 

\begin{table}
	\caption{Forms of upward extensions (Biu-Mandara)}
	\label{tab:uplang}
	\begin{tabular}{lllll}
		\lsptoprule
		&     Forms    &  Subbranch & Group & Subgroup        \\
		\midrule
		Psikye&	\textit{(a)me, ate}	& North & Higic & \\
		Lamang&	\textit{f(i)} &	North & Lamang-Hdi & \\
		Dghwede&	\textit{gwá, fè} &	North & M.-M.-M. & Mandaraic \\
		Parkwa&	\textit{u, əlu, alu} &	North & M.-M.-M. & Mandaraic  \\
		Glavda&	\textit{(i)t} &	North & M.-M.-M. & Mandaraic  \\
		Malgwa&	\textit{tə́} &	North & M.-M.-M. & Mandaraic  \\
		Gude&	\textit{gi} &	South & Bataic & \\
		Sakun&	\textit{rá, má} &	South & Sukur & \\
%		Masana&	\textit{kúló}	& Masa North & &  \\
		\lspbottomrule
	\end{tabular}
\end{table}

\begin{table}
	\caption{Forms of downward extensions (Biu-Mandara)}
	\label{tab:downlang}
	\begin{tabular}{lllll}
		\lsptoprule
		&     Forms    &  Subbranch & Group  & Subgroup       \\
		\midrule
		Psikye	& \textit{awa, (i)yi}	&	North	&	Higic &	\\
		Lamang	& \textit{gaa, ɗi}	&	North	&	Lamang-Hdi &	\\	
		Bura	& \textit{hi}	&	North		&	M.-M.-M. & Bura-Marghi			\\
		Margi & 	\textit{ía} & 		North		&	M.-M.-M. & Bura-Marghi			\\
		Huba	& \textit{yà}	& 	North		&	M.-M.-M. & Bura-Marghi			\\
		Dghwede&	\textit{tígè, àwáyà, àwí}	& BM	North	&	M.-M.-M. & Mandaraic	\\
		Parkwa&	\textit{aha, əla, ada}	& North	&	M.-M.-M. & Mandaraic  \\
		Glavda&	\textit{(x)i}	&	North	&	M.-M.-M.	& Mandaraic \\	
		Mada 	& \textit{ra}	&	North	&	M.-M.-M. & Mofuic	\\	
		Gude&	\textit{gərə, paa}	&	BM	&	Bataic &	\\	
		Sakun&	\textit{xá} &	South	&	Sukur	& \\
		Ga'anda &	\textit{xa}	&	South & Teraic & \\
%		Masana&	\textit{gà} &	Masa North	&	& \\		
		\lspbottomrule
	\end{tabular}
\end{table}

\newpage
\section{Other directional extensions} \label{sec:diachronyother}

The forms of the nine `onto' directional extensions are listed in Table \ref{tab:ontolang}. There are two clear sets of cognates. The Bura-Margi languages have an \textsc{onto} extension similar to \textit{nger}. \citet[292]{hoffmannUntersuchungenZurStruktur1955} notes that this extension is similar in form to the noun \textit{kir} `head' -- a path of grammaticalization that is well attested outside of Chadic languages \citep[169--170]{heineWorldLexiconGrammaticalization2002}. Three of the four \textsc{onto} extensions in Mandaraic languages have the form \textit{ar(ə)}. This could potentially be a shortened cognate form of the Bura-Margi `onto' extensions.

\begin{table}
	\caption{Biu-Mandara North languages with an ONTO extension}
	\label{tab:ontolang}
	\begin{tabular}{llll}
		\lsptoprule
		&     Form    &  Group & Subgroup         \\
		\midrule
		Buduma& \textit{gə̀rə́} & Kotoko-Buduma  & \\
		Bura& \textit{nkir} & M.-M.-M. & Bura-Margi \\
		Margi& \textit{ŋgéri} & M.-M.-M. & Bura-Margi \\
		Huba& \textit{ngərì} & M.-M.-M. & Bura-Margi \\
		Dghwede& \textit{gá} & M.-M.-M. & Mandaraic \\
		Parkwa& \textit{arə} & M.-M.-M. & Mandaraic \\
		Malgwa& \textit{ar} & M.-M.-M. & Mandaraic \\
		Wandala& \textit{ar} & M.-M.-M. & Mandaraic \\
		Mada& \textit{ka}, \textit{fa} & M.-M.-M. & Mofuic \\
		\lspbottomrule
	\end{tabular}
\end{table}

As already noted in Chapter \ref{ch:distribution}, Section \ref{sec:onto2}, the `under' directional extension has a form featuring \textit{s} in the three Mandaraic languages (Biu-Mandara North) where it appears. The source of these likely cognates is not immediately apparent. It is probably only a coincidence that the ventive directional verb roots in Mandaraic languages are all of the form \textit{s} or similar as well.\footnote{In addition, \citet[3]{blenchBuraVerbalExtensions2010} lists two forms of \hyperref[bura]{Bura} extensions meaning `under'. One is only shown combined with the verb \textit{lə} `go': \textit{ləkəra} `to go under'. The other is in a list of ``adpositional suffixes'' meaning ``under, below, beneath'' but without any example sentences. Due to the lack of examples, Bura has not been included as a language with an `under' directional extension.}

In regard to homeward extensions, there is good evidence that these are derived from a noun. The homeward directional suffix \textit{-gha} in \hyperref[lama]{Lamang} is said to be derived from the noun \textit{tə́ghà} `compound, residence, home' \citep[140]{wolffLamangLanguageDictionary2015}. This can be compared with \hyperref[psik]{Psikye}, where there is phonological evidence that the word \textit{gɛ} `compound, home' may be grammaticalizing since, when it occurs next to a basic motion verb, it can trigger palatalization of the verb, as seen in (\ref{ex:psik14}) (<š> = [ʃ]). However, there is no evidence that this form has generalized to other verbs as part of the verbal morphology. 

\ea\label{ex:psik14}
\langinfo{\hyperref[psik]{Psikye} verb \textit{sé} `come' phonologically attached to \textit{gɛ} `home'}{}{\cite[118]{smithKapsikiLanguage1969}} \\
\gll ša=gɛ ŋkɛ́ le yɛmú \\
come=home \gl{poss}.\gl{3}\gl{sg}.\gl{f} with water	\\
\glt `She is coming home with the water.'
\z 

The other language with a homeward directional extension is \hyperref[bura]{Bura}. In Bura, the homeward suffix \textit{-vi} is compared to a word \textit{avi} `in the house, at home' \citep[301]{hoffmannUntersuchungenZurStruktur1955} or to a noun \textit{vi} `place' \citep[7]{blenchBuraVerbalExtensions2010}. In the same subgroup, \citet[24]{muazuModernKilbaDictionary2015} list the \hyperref[huba]{Huba} verb \textit{gúví} `to go into a house or room' which seems to show a parallel with the homeward suffixes found in other Bura-Marghi languages, but in Huba it is not analyzed as a suffix. 

The one language that is said to have a suffix with a possible bushward directional interpretation is \hyperref[psik]{Psikye}. A close parallel to the \hyperref[psik]{Psikye} suffix is found in \hyperref[saku]{Sakun}. \citet[122]{thomasGrammarSakunSukur2014} reports: ``The directional extension \textit{-ᵐta} `to the bush' only combines with the root \textit{dza} ‘to go’ and its perfective form \textit{ra}.'' Since this form is only found with one verb, it is not considered part of the verbal morphology. \citet[163]{smithKapsikiLanguage1969} notes that the Psikye noun \textit{gambá} `bush' is a less likely source for the suffix than a prepositional phrase \textit{kwa mté} `in the bush' which is analyzed as containing the verbal stem \textit{mté} `die'. A similar situation is seen in Sakun where the noun is \textit{gaᵐbá} `bush' but the data contains another expression \textit{xaᵐtá} `bush' as well. 

Possible cognates of bushward suffixes are the itive suffix \textit{-mta} in \hyperref[bura]{Bura} \citep{hoffmannUntersuchungenZurStruktur1955,blenchBuraVerbalExtensions2010} and the \hyperref[ciba]{Cibak} suffix \textit{-nta} which may be cognate even though it is not attested with a directional function \citep[134]{hoffmannZurSpracheCibak1955}. On the connection between bushward and itive direction, there is a semantic (though not syntactic) parallel found in the East Chadic language \hyperref[lele]{Lele}. The locative adverbial \textit{cáaní} ‘in/at the bush’, as in (\ref{ex:lele7}), also has an itive directional function, as in (\ref{ex:lele0}), and a resultative itive function, as in (\ref{ex:lele9}).

\ea\label{ex:lele7}
\langinfo{\hyperref[lele]{Lele} adverbial \textit{cáaní} ‘in the bush’ }{}{\cite[211]{frajzyngierGrammarLele2001}} \\
\gll màdú	nè	ná	cáaní	lay \\
vitex	\gl{cop}	\gl{asoc}	field	also \\
\glt `The fruits of vitex are in the bush.'

\ex\label{ex:lele0}
\langinfo{\hyperref[lele]{Lele} verbs \textit{gìr è} `run go' with adverbial \textit{cáaní} }{}{\cite[184]{frajzyngierGrammarLele2001}} \\
\gll gú∫yé ɗeŋlí se jìb se gìr è cáani	\\
spider hear \gl{incept} push \gl{incept} run go bush\\
\glt `Spider heard it, got up, and ran away.'

\ex\label{ex:lele9}
\langinfo{\hyperref[lele]{Lele} verb \textit{sáde} ‘clean’ with adverbial \textit{cáaní}}{}{\cite[198]{frajzyngierGrammarLele2001}} \\
\gll ŋ	sáde	kùbàrò	cáaní \\
\gl{1}\gl{sg}	clean.\gl{fut}	blood	\gl{itive} \\
\glt `I will clean the blood away.'
\z 

Finally, in one language, \hyperref[bura]{Bura}, there are marginal examples of a directional extension meaning `toward the side, or next to' \citep{blenchBuraVerbalExtensions2010}. This suffix is attested with at least one intransitive motion verb, shown in (\ref{ex:bura12}), as well as transitive motion verbs, as in (\ref{ex:bura13}). \citet[275]{hoffmannUntersuchungenZurStruktur1955} also discusses the \textit{-dza} suffix but describes its meaning as more generally locative, etymologically related to a preposition \textit{adza} `to, at'. Hoffman raises doubts about the analysis of \textit{-dza} as a suffix and explicitly says that the form \textit{lidza}, as in (\ref{ex:bura12}), is not a suffix, but merely a phonetic amalgamation of the verb and preposition. 

\ea\label{ex:bura12}
\langinfo{\hyperref[bura]{Bura} verb \textit{li} `go' with the `beside' suffix}{}{\cite[4]{blenchBuraVerbalExtensions2010}} \\
\gll li-dza \\
go-\textsc{side} \\
\glt `to go close to'

\ex\label{ex:bura13}
\langinfo{\hyperref[bura]{Bura} verb \textit{bwal} `move' with the `beside' suffix}{}{\cite[4]{blenchBuraVerbalExtensions2010}} \\
\gll bwal-dza \\
move-\textsc{side} \\
\glt `to place an object beside another'
\z 

\section{Conclusion}

In a review of an etymological dictionary, \citet{takacsReviewArticlesFirst2017} comments that, in regard to the task of historical reconstruction, ``the quantity and the diversity of the Chadic daughter languages is all too enormous to overcome.'' While this pessimistic sentiment is relatable when considering the diachrony of directional extensions, there are some plausible generalizations that can be made about the sources of these morphemes. 

For West Chadic languages, earlier hypotheses regarding the reconstruction of ventive extensions seem to be on the right track, and there is a plausible scenario in which these ventive extensions were formed from multiverb constructions \citep{frajzyngierVentiveCentrifugalChadic1987}. It seems less likely that any ventive directional extensions can be reconstructed for Proto-Chadic. The origins of ventive extensions in Biu-Mandara, Masa and East Chadic languages are not certain, though there are some suggestions of adverbial or prepositional sources. 

Since the widest range of directional extensions only occur in Biu-Mandara languages, with no evidence of traces in other branches, it can be assumed that the vertical and boundary-crossing directional extensions are the result of innovation after the split of Biu-Mandara languages from other branches. The variety of forms for these extensions suggests multiple different sources which could be verbal, adverbial, prepositional or nominal. \citet[16]{blenchBuraVerbalExtensions2010} hypothesizes that this diversity can be explained by assuming that ``the concept of a complex semantics using verbal extensions is spread through metatypy, the diffusion of a strategy from language to language, rather than direct borrowing.'' The data on co-occurrence of directional extensions (Chapter \ref{ch:distribution}, Section \ref{sec:co-occurrence}) suggest that the less common `onto' and `under' directional extensions, as well as bushward and homeward extensions, only develop in languages that already have vertical and boundary-crossing directional extensions. 

Given that directional extensions commonly express other meanings, including tense-aspect, it is plausible that directional extensions in at least some languages have grammaticalized away from their directional function to only express a more abstract aspectual meaning. In order to investigate this possibility, there would need to be a systematic comparison of tense-aspect marking across Chadic languages. This would ideally be complemented by continued work on understand sound correspondences, especially in West Chadic, East Chadic and Masa. As \citet{souagEkkehardWolffHistorical2023} states in a review of \citet{wolffHistoricalPhonologyCentral2021}: ``Better documentation of these languages is thus all the more desirable for the purposes of historical linguistics, as well as in its own right.'' Morphosyntactic research, like the study of directionals, creates a valuable vantage point for comparing Chadic languages and reconstructing their histories. The data presented in this chapter has provided new insights and also raised new questions as an opportunity to develop a deeper understanding of the relationships between Chadic languages and beyond.
